\section{Conclusion}

There is a growing demand for real-time interaction with agentic AI technologies to align their behavior with user demands; however, complex agentic workflows typically add several seconds of latency due to the time for the model to reason for each action as well as long tool latencies.
In this work, we present an algorithm and system design which enables real-time interaction with agentic systems with low latency.
We propose \textit{Asynchronous I/O}, where we decouple the agent's thinking and acting from delays with information arriving from the user and environment.
Our \textit{Speculative Tool Calling} approach allows the model to launch tool calls and later correct any mistakes, and ensures that sensitive tools do not execute until we have ensured their arguments are correct based on the full user input.
We present an event-driven inference framework which contains a task execution manager as well as a core agent loop which supports interrupting agent execution when new information arrives.
We also propose a clock-based training methodology for generating synthetic data for SFT.
Altogether, our \textit{Speculative Interaction Agent} methodology provides up to 2.2$\times$ speedup relative to the reason and act baseline.


\section{Acknowledgements}

We acknowledge gracious support from the Furiosa AI, Intel, Apple, NVIDIA, Macronix, Panasonic, and Mozilla teams. 
Furthermore, we appreciate support from Google Cloud, the Google TRC team, and specifically Jonathan Caton, Divvy Thakkar, and Prof. David Patterson.
Prof. Keutzer’s lab is also sponsored by funding through BDD and BAIR. 
We also acknowledge support by the Director, Office of Science, Office of Advanced Scientific Computing Research, of the U.S. Department of Energy under
Contract No. DE-AC02-05CH11231. Michael W. Mahoney acknowledges
DARPA, NSF, the DOE Competitive Portfolios grant, and
the DOE SciGPT grant. Our conclusions do not necessarily
reflect the position or the policy of our sponsors, and no
official endorsement should be inferred.